% FILE: 03_architecture_and_primitives.tex
% PROJECT: ma — M&A Claim Taxonomy and Board Projection Surface

\section{Architecture and Primitives}
\label{sec:ma-architecture}

\subsection{Core Objects}
\label{sec:ma-arch-objects}

The \texttt{ma/} framework is built on six core object types, each defined in a dedicated
specification file and referenced by the master consolidation document.

\paragraph{Claim.} A typed assertion about an organization's process, expressed as a
mathematical formula whose variables are derivable from an XES or OCEL 2.0 event log. The
nine claim taxonomies enumerate specific claim classes with their formulas:

\begin{itemize}
  \item \textbf{C-EBITDA-001} --- EBITDA Optimization:
    $E = V_{\text{annual}} \times (r_{\text{baseline}} - r_{\text{target}}) \times
    \bar{C}_{\text{rework}}$
  \item \textbf{C-WC-002} --- Working Capital Release:
    $WC = (\text{Revenue}_{\text{credit\_annual}} / 365) \times (T_{\text{AR, baseline}} -
    T_{\text{AR, target}})$
  \item \textbf{C-RISK-003} --- SLA Penalty Exposure:
    $L_{\text{SLA}} = \sum_{c \in C_{\text{late}}} P_{\text{penalty}}(c) +
    \sum_{c \in C_{\text{active}}} \Pr(T(c) > T_{\text{SLA}} \mid \sigma_c) \times
    P_{\text{penalty}}(c)$
  \item \textbf{C-RISK-004} --- Compliance Leakage:
    $L_{\text{compliance}} = \sum_{r \in \mathcal{R}} (N_{\text{violations}}(r, L) \times
    F_{\text{statutory}}(r) + \Pr(\text{Audit}_{\text{ext}}) \times F_{\text{systemic}}(r))$
  \item \textbf{C-RESIDUAL-005} --- Process Standardization: residual weight
    $W_R = |R| / |L|$ and residual entropy $H_R = -\sum_{v \in V_R} P(v) \log_2 P(v)$
\end{itemize}

\paragraph{Receipt.} A JSON document conforming to RFC 8785 (JCS canonical form), containing
the slide UUID, BLAKE3 hashes of event log and process model, wasm4pm query definition,
numeric verification results, and an Ed25519 validator signature. Five materialized receipts
exist in \texttt{receipts/}; additional receipts are specified by schema.

\paragraph{Event Log.} An XES (IEEE 1849-2016) or OCEL 2.0 event log extracted from a source
ERP or CRM system, with W3C PROV-O provenance lineage mapping events to source database WAL
sequence numbers. Each trace $\sigma = \langle e_1, e_2, \ldots, e_n \rangle$ forms a hash
chain: $\mathcal{H}(e_j) = \text{BLAKE3}(e_j \Vert \mathcal{H}(e_{j-1}) \Vert
\text{Sig}_{\text{system}}(e_j))$.

\paragraph{Process Model.} A Workflow Net, POWL structure, or process tree discovered from the
event log via Inductive Miner (Leemans 2013) or equivalent algorithm. Models must satisfy
soundness: option to complete ($\forall M \in [M_0\rangle$, $[o] \in [M\rangle$), liveness
($\forall t \in T$, $\exists M', M'' : M \xrightarrow{*} M' \xrightarrow{t} M''$), and
boundedness ($\exists k \in \mathbb{N}^+ : \forall M \in [M_0\rangle$, $\forall p \in P$,
$M(p) \le k$).

\paragraph{Virtual Data Room (VDR).} A structured evidence archive containing six directories:
\texttt{/event-logs/}, \texttt{/models/}, \texttt{/receipts/}, \texttt{/conformance/},
\texttt{/provenance/}, \texttt{/synergy/}. The VDR is the delivery vehicle by which the seller
provides all evidence necessary for independent buyer replay.

\paragraph{Board Presentation.} A PowerPoint deck produced by the seven-step Board Projection
Renderer pipeline. Every slide assertion is sealed to a receipt via the Slide-to-Receipt Map.

\subsection{Admissible Transitions}
\label{sec:ma-arch-transitions}

The framework defines a strict evidence-escalation protocol. A claim advances through the
Board Claim Taxonomy tiers only when the evidence minimum for the target tier is met:

\begin{itemize}
  \item T1 $\to$ T2: Volume metrics and coverage mapping available from system-of-record.
  \item T2 $\to$ T3: OCEL/XES event log available; process model discovered; conformance
    report with $f \ge 0.95$, $p \ge 0.90$ produced.
  \item T3 $\to$ T4: Event timestamps admitted; case duration CDF computed; performance
    claims numerically bounded.
  \item T4 $\to$ T5: \texttt{wasm4pm} running on live data; refusal evidence---queries that
    correctly reject non-conforming traces---demonstrated.
  \item T5 $\to$ T6: Process variant analysis performed; conformance delta quantified;
    remediation receipt produced with root-cause trace.
\end{itemize}

The Four-Pillar Board Admissibility Protocol governs admission at T3 and above: Pillar A
(Event Log Integrity), Pillar B (Mathematical Conformance Bounds), Pillar C (Structural Model
Soundness), Pillar D (Cryptographic Slide-to-Receipt Mapping) must all pass.

\subsection{Boundaries}
\label{sec:ma-arch-boundaries}

The \texttt{ma/} module is bounded on three sides:

\textbf{Upstream boundary:} \texttt{wasm4pm} is the execution engine. Conformance queries
must be expressible as \texttt{wasm4pm} query URIs. The module does not define the query
execution engine---it defines the claim taxonomy, admissibility protocol, and receipt schema
that constrain what queries must produce.

\textbf{Downstream boundary:} The Board Projection Renderer (Rust: \texttt{Tera2},
\texttt{pptx-rs}, \texttt{calamine}) is the rendering surface. The module does not define
rendering implementation---it defines the claims JSON contract that the renderer must consume.

\textbf{Lateral boundary:} The five-step Auditor Evidence Path defines the external-auditor
interface. Auditors receive the VDR and execute the protocol; they do not modify receipts or
re-run discovery. Replay tolerance is $\Delta f, \Delta p < 10^{-6}$ across independent runs.

\subsection{Failure Modes Refused}
\label{sec:ma-arch-failure-modes}

The framework explicitly refuses four failure modes that are endemic to traditional M\&A
diligence:

\begin{enumerate}
  \item \textbf{Narration without replay:} A claim that cannot be independently replayed
    from VDR contents is not a board-admissible claim---it is narration. The Slide-to-Receipt
    Map makes orphaned assertions structurally impossible: any slide without a receipt UUID
    fails the Pillar D check.
  \item \textbf{Statistical sampling without event log:} T2 coverage claims are explicitly
    lower-tier. A process with 100\% coverage and 40\% conformance fitness is a liability, not
    an asset. The taxonomy requires escalation to T3 before operational claims are priced.
  \item \textbf{Synergy claims without behavioral similarity proof:} Post-merger synergy
    claims require Behavioral Profile Similarity $\text{Sim}(M_1, M_2, C) \ge 0.75$ (Weidlich
    2011) and Execution Path Intersection evidence. A synergy claim without bisimulation
    equivalence evidence is inadmissible under v30.1.1.
  \item \textbf{Soundness-uncertified process models:} All Workflow Nets submitted as
    conformance evidence must carry a soundness certificate. An unsound model can admit traces
    that replay against it without meaningful conformance---a form of false positive that inflates
    fitness scores. The module requires $k$-boundedness $k \le 5$ for financial processes.
\end{enumerate}
